\documentclass[12pt,letterpaper]{article}
\usepackage[utf8]{inputenc}
\usepackage[spanish]{babel}
\usepackage{amsmath}
\usepackage{amsfonts}
\usepackage{amssymb}
\usepackage[margin=2.5cm]{geometry}

\begin{document}

\begin{center}
    \Large \textbf{Evaluación de Examen 2: Unidad II} \\
    \large Física para Ciencias de la Salud (005-1713) \\
    \vspace{0.5cm}
    \textbf{Estudiante:} Michelle Gabriela López \quad \textbf{C.I.:} 33.483.425
    \vspace{0.3cm} \\
    \large \textbf{Calificación Final: 10/10 puntos}
\end{center}

\vspace{0.5cm}

\section*{Análisis de Resultados}

\subsection*{1. Cinemática (Aceleración media)}
\textbf{Procedimiento y resultado:} Extrae formalmente los datos incluyendo $\Delta t = 0,15 \text{ s}$. Plantea la ecuación de aceleración media y realiza la sustitución correcta restando las velocidades ($0,6 \text{ m/s} - 0 \text{ m/s}$). El cálculo aritmético arroja el resultado analítico exacto de $4 \text{ m/s}^2$. \\
\textbf{Veredicto:} \textbf{Correcto} (2/2 pts).

\subsection*{2. Dinámica (Leyes de Newton)}
\textbf{Procedimiento y resultado:} Muestra un despeje impecable a partir de la Segunda Ley de Newton ($a = \frac{F}{m}$). Sustituye los valores dividiendo $150 \text{ N}$ entre $25 \text{ kg}$, logrando un resultado perfecto de $6 \text{ m/s}^2$. Destaca de forma excelente al mostrar a un lado el análisis dimensional ($\frac{\text{kg} \cdot \text{m/s}^2}{\text{kg}} = \text{m/s}^2$) y aplicar notación científica ($6,0 \times 10^0 \text{ m/s}^2$). \\
\textbf{Veredicto:} \textbf{Correcto} (2/2 pts).

\subsection*{3. Trabajo Mecánico}
\textbf{Procedimiento y resultado:} Identifica el ángulo como $0^\circ$ de antemano y plantea la fórmula del trabajo ($W = F \cdot d$). Efectúa correctamente la multiplicación directa ($400 \text{ N} \cdot 0,8 \text{ m}$). El cálculo da exactamente $320 \text{ J}$, y reporta su notación científica correcta como $3,20 \times 10^2 \text{ J}$. \\
\textbf{Veredicto:} \textbf{Correcto} (2/2 pts).

\subsection*{4. Energía Potencial}
\textbf{Procedimiento y resultado:} Plantea el producto de las tres variables según la ecuación de la energía potencial gravitatoria ($E_p = m \cdot g \cdot h$). Ejecuta la multiplicación en pasos separados ($1,2 \times 9,8 = 11,76$, luego multiplica por $1,5$), obteniendo el resultado verificado de $17,64 \text{ J}$. La conversión a notación científica ($1,76 \times 10^1 \text{ J}$) es válida. \\
\textbf{Veredicto:} \textbf{Correcto} (2/2 pts).

\subsection*{5. Potencia Mecánica}
\textbf{Procedimiento y resultado:} Relaciona el trabajo efectuado y el tiempo planteando la fracción respectiva ($P = \frac{W}{t}$). El desarrollo de la división de $75 \text{ J} / 60 \text{ s}$ arroja la cifra correcta de $1,25 \text{ W}$, indicando correctamente su notación en $1,25 \times 10^0 \text{ W}$. \\
\textbf{Veredicto:} \textbf{Correcto} (2/2 pts).

\vspace{0.5cm}
\section*{Comentarios Generales y Sugerencias}
¡Felicidades por obtener una calificación perfecta! El examen demuestra un nivel de excelencia total, no solo en la asimilación de los conceptos teóricos, sino en el manejo avanzado de las matemáticas y la notación técnica.
\begin{itemize}
    \item El hábito de expresar los resultados finales acompañados de su respectiva \textbf{Notación Científica} es sumamente valorado. Evidencia un dominio destacable sobre las leyes de la potenciación (como saber que $10^0 = 1$ para valores de una sola cifra).
    \item Mostrar de manera explícita el **análisis dimensional** (como el desglose de los Newtons en la Pregunta 2 para cancelar kilogramos) es la práctica ideal que asegura cálculos correctos en investigaciones científicas futuras.
    \item ¡Sigue con este mismo entusiasmo y excepcional nivel de excelencia académica!
\end{itemize}

\end{document}